\chapter{Defending this work}\label{ch:defence}

These are the questions most likely to come up, with honest answers.

\begin{careful}[Written on 5 November]
This chapter predates MAGIC-UNet. Its recommendation of FISTA-Net is now superseded on COMSOL score: MAGIC-UNet run 2
scores 27.0\% without being told the standoff. Chapter~\ref{ch:magic_results} has the new questions and answers, and the
limits of that result.
\end{careful}

\section{About the source}

\textbf{What is your network, and where does it come from?}\\
Eq.~20 of Liu, Chen, Wang and Yin, ICLR 2019: convolutional LISTA. The extension to deblurring,
$d_m=h*g_m$, is ours, and no published result covers it.

\textbf{Is it ALISTA?}\\
No. ALISTA computes its weights analytically. We learn ours. Eq.~26 was tried and dropped.

\textbf{Which network do you recommend?}\\
FISTA-Net run 2, from Xiang, Dong and Yang. It has the lowest simulated validation error among runs without our floor and loss
changes, and it scores 33.3\% on COMSOL. Strict Eq.~20 at 33.8\% is effectively tied.

\textbf{Which parts did you invent?}\\
For strict runs: the blurred dictionary, the padding and mask, and the dictionary learning method. Chapter~\ref{ch:differences}
lists all training details. The floor and the relative loss were additions, and strict runs removed them.

\section{About the method}

\textbf{Why not just use FISTA, since it wins?}\\
On accuracy alone, you should. FISTA-300 scores 30.5\%, against 33.3\% for the selected network. Strict Eq.~20 wins at equal
depth, 33.8 against 37.9\% at 64, but only on COMSOL, not on simulated validation. Run time has not been measured, so I
cannot claim any network is faster.

\textbf{Was FISTA tuned fairly?}\\
Yes. Its $\lambda$ came from eight values on 16 simulated crops, and the best one lies inside the grid. COMSOL was not used. The
search was widened after the network scores were known, and you should say so.

\textbf{Did FISTA-Net fail?}\\
Run 1 scored 39.4\%, but it was stopped early and its start values gave a gradient step 17 times weaker than the authors' code.
Run 2 fixed both and scored 33.3\%, so the low score came from our setup. Chapter~\ref{ch:fistanet} has the numbers.

\textbf{Is your network doing sparse coding?}\\
Structurally yes. In behaviour, mostly no. Strict codes have high density, and some thresholds are zero, so it acts
like a learned deconvolution.

\textbf{Does the ALISTA convergence theorem apply?}\\
No. It assumes much sparser codes than ours.

\textbf{Why does $h$ stay fixed?}\\
Because it is physics: upward continuation in a source-free region. We know it exactly, so learning it would only
waste data. The empirical $k_0$ was removed because it breaks harmonicity.

\section{About the data}

\textbf{Why train on simulation?}\\
Only 15 real pairs from two wire shapes exist, far too few for millions of parameters. Simulation derives all
components from one current map, so it is physically consistent.

\textbf{Is the test set really clean?}\\
Models and checkpoints were chosen without COMSOL. But the generator's parameter ranges were designed from these coils,
and we have now compared many methods on them. So it is clean of reconstruction scores, but not completely independent.

\textbf{Why 33.8\% if the method is good? That sounds bad.}\\
Because of how the score works. COMSOL's speckle sets a floor around 15--19\%, the blurry input itself is near 77\% on a
signal-rich crop, and the Wiener filter gets 49.3\%.

\section{About the results}

\textbf{Why did simulated validation disagree with COMSOL?}\\
The simulator differs from COMSOL, for example in its speckle and peaks, and the difference is not fully understood.
The sparse variant won on simulation and lost on COMSOL, so simulated gains do not always transfer.

\textbf{How confident are you in a one-point difference?}\\
Not very. There is one seed per run, and no error bars.

\section{What is not yet known}

\begin{itemize}
\item Where FISTA-Net would settle with a decaying learning rate, or more updates.
\item How stable FISTA-Net's score is across seeds. Its validation error is still noisy.
\item Whether the LFISTA-style network improves at $\lambda=0.009$, or with a higher limit on $\beta$.
\item Actual run times. Without them, no speed claim is possible.
\item Why strict Eq.~20's gain over FISTA at equal depth appears on COMSOL and not on simulated validation.
\item Whether the early validation peak in MSE runs is overfitting or a loss mismatch.
\item Why COMSOL decays faster than free-space physics.
\item Whether results hold on a third wire shape.
\end{itemize}
